\chapter*{Определения, обозначения и сокращения}
\addcontentsline{toc}{chapter}{Определения, обозначения и сокращения}

\newcommand{\glossitem}[2]{\textbf{\textit{#1}} - #2\par}

\glossitem{Биоинформатика}{область, связанная с вычислительной обработкой, анализом и интерпретацией биологических данных.}

\glossitem{Биологическая последовательность}{строковое представление DNA, RNA или белка, состоящее из символов соответствующего биологического алфавита.}

\glossitem{Сонификация (Sonification)}{преобразование данных в звуковую форму для восприятия или анализа.}

\glossitem{Биосонификация (Biosonification)}{преобразование биологических данных в звуковую или музыкальную форму для восприятия человеком.}

\glossitem{Символическая музыка (Symbolic music)}{представление музыки в виде дискретных событий, например нот, длительностей, аккордов и временных позиций, а не в виде аудиосигнала.}

\glossitem{MusicLang}{высокоуровневый текстовый язык описания музыки через аккорды, тональности и относительные ноты.}

\glossitem{MIDI (Musical Instrument Digital Interface)}{стандарт цифрового представления музыки в виде событий, таких как note on, note off, pitch и velocity.}

\glossitem{Модель-оценщик}{модель или алгоритмический контур, используемый для оценки качества сгенерированного музыкального фрагмента.}

\glossitem{Генеративно-состязательная сеть (GAN, Generative Adversarial Network)}{архитектура из двух сетей, генератора и дискриминатора, обучающихся в состязательном режиме.}

\glossitem{Трансформер (Transformer)}{архитектура нейронной сети, основанная на механизме внимания и не использующая рекуррентные связи.}

\glossitem{Baseline}{базовый вариант сравнения, относительно которого оценивается качество разработанной модели.}

\glossitem{POP909}{датасет из 909 популярных песен в MIDI-формате с аннотациями аккордов и мелодий.}

\glossitem{MAESTRO}{датасет фортепианных MIDI-записей, часто используемый в исследованиях генерации символической музыки.}

\glossitem{FASTA}{текстовый формат для представления биологических последовательностей DNA, RNA и protein.}

\glossitem{DNA}{дезоксирибонуклеиновая кислота; в работе рассматривается как последовательность нуклеотидных символов.}

\glossitem{RNA}{рибонуклеиновая кислота; в работе рассматривается как последовательность нуклеотидных символов с урацилом вместо тимина.}

\glossitem{Protein}{белковая последовательность, представленная символами аминокислот.}

\glossitem{Авторегрессионная модель (Autoregressive model)}{модель, которая генерирует последовательность по шагам, предсказывая следующий токен на основе уже сгенерированного контекста.}

\glossitem{Функция потерь (Loss)}{численная функция ошибки, оптимизируемая при обучении модели.}

\glossitem{PyTorch}{фреймворк глубокого обучения с динамическими вычислительными графами.}

\glossitem{Checkpoint}{сохраненное состояние обученной модели, которое можно загрузить для продолжения обучения или генерации.}

\glossitem{RefSeq}{база эталонных биологических последовательностей, используемая как источник геномных данных.}

\glossitem{Нуклеотид (Nucleotide)}{структурная единица DNA или RNA, состоящая из азотистого основания, сахара и фосфатной группы.}

\glossitem{GC-content}{доля гуанина и цитозина в нуклеотидной последовательности.}

\glossitem{k-mer}{подпоследовательность длины k, например триплет нуклеотидов при k=3.}

\glossitem{Энтропия последовательности (Sequence entropy)}{мера неопределенности или разнообразия символов в последовательности.}

\glossitem{Кодон (Codon)}{тройка нуклеотидов, кодирующая одну аминокислоту.}

\glossitem{Открытая рамка считывания (ORF, Open Reading Frame)}{участок последовательности между стартовым и стоп-кодонами.}

\glossitem{GRAVY (Grand Average of Hydropathy)}{средняя гидрофобность аминокислот в белковой последовательности.}

\glossitem{Изоэлектрическая точка (pI, Isoelectric point)}{pH, при котором белок имеет нулевой суммарный заряд.}

\glossitem{Embedding}{численное векторное представление объекта, например биологического фрагмента или токена, пригодное для обработки нейросетью.}

\glossitem{Control profile}{компактный вектор управляющих признаков, описывающий целевые музыкальные свойства фрагмента: темп, плотность нот, частоту смены гармонии, регистр, сложность и модальность.}

\glossitem{Управляющий сигнал (conditioning signal)}{входная информация, задающая условия генерации модели и влияющая на создаваемый результат; в данной работе таким сигналом выступают биологические признаки и полученные из них векторные представления.}

\glossitem{Bio conditioning}{использование биологического вектора или профиля признаков как управляющего условия для нейросетевой генерации.}

\glossitem{Bio->Harmony->Melody}{иерархическая схема генерации, в которой биологические признаки сначала управляют построением гармонии, а затем мелодия генерируется с учетом этой гармонии.}

\glossitem{Гармония (Harmony)}{аккордовая структура музыкальной композиции.}

\glossitem{Аккорд (Chord)}{одновременное звучание трех и более нот, образующих гармоническую структуру.}

\glossitem{Мелодия (Melody)}{последовательность нот, образующая музыкальную линию.}

\glossitem{Такт (Bar, Measure)}{единица музыкального времени, содержащая фиксированное количество долей.}

\glossitem{Октава (Octave)}{интервал между двумя нотами, частота которых отличается в два раза.}

\glossitem{Полутон (Semitone)}{минимальный интервал в западной музыкальной системе, равный одной двенадцатой октавы.}

\glossitem{Токенизация (Tokenization)}{преобразование данных в последовательность дискретных токенов для обработки моделью.}

\glossitem{Memory tokens}{служебные токены, через которые биологический embedding передается трансформеру как дополнительный контекст.}

\glossitem{Top-k matching}{подбор одного из k ближайших музыкальных сегментов по расстоянию между биологическим управляющим профилем и музыкальными дескрипторами.}

\glossitem{Train/validation/test split}{разбиение данных на обучающую, валидационную и тестовую части для обучения, настройки и независимой проверки модели.}

\glossitem{Epoch}{один полный проход алгоритма обучения по обучающей выборке.}

\glossitem{Gradient accumulation}{техника обучения, позволяющая эмулировать большой batch size на ограниченной памяти.}

\glossitem{Mixed precision}{техника обучения нейронных сетей с использованием 16-битных и 32-битных чисел для экономии памяти.}

\glossitem{Cross-entropy loss}{функция потерь для задачи предсказания следующего токена; показывает, насколько вероятностное распределение модели отличается от правильного ответа.}

\glossitem{Validation loss}{значение функции потерь на валидационной выборке, которая не используется для обновления параметров модели.}

\glossitem{Valid MIDI rate}{доля сгенерированных MIDI-файлов, которые корректно открываются и содержат ожидаемую музыкальную структуру.}

\glossitem{Chord-tone ratio}{доля мелодических нот, принадлежащих текущему аккорду; используется как метрика согласованности мелодии и гармонии.}

\glossitem{Self-similarity}{метрика повторяемости или внутренней согласованности музыкального фрагмента.}

\glossitem{Абляционное исследование (Ablation study)}{эксперимент, в котором отдельные компоненты модели или входных данных отключаются или заменяются, чтобы оценить их вклад в результат.}

\glossitem{Real bio}{вариант абляционного сравнения, в котором модель получает настоящий биологический embedding и настоящий управляющий профиль.}

\glossitem{Shuffled bio}{вариант абляционного сравнения, в котором биологическое условие перемешивается между фрагментами для проверки вклада связи ``биоданные - музыка''.}

\glossitem{Neutral bio}{вариант абляционного сравнения, в котором биологическое условие заменяется нейтральным или усредненным представлением.}

\glossitem{Seed}{фиксированное начальное значение генератора случайных чисел, обеспечивающее воспроизводимость эксперимента.}

\glossitem{Автоэнкодер (Autoencoder)}{нейронная сеть, обучающаяся сжимать данные в латентное представление и восстанавливать их обратно.}

\glossitem{Вариационный автоэнкодер (VAE, Variational Autoencoder)}{генеративная модель, использующая вероятностное латентное пространство.}

\glossitem{Длинная краткосрочная память (LSTM, Long Short-Term Memory)}{тип рекуррентной нейронной сети с механизмом управления памятью.}

\glossitem{Рекуррентная нейронная сеть (RNN, Recurrent Neural Network)}{нейронная сеть, обрабатывающая последовательности с использованием внутреннего состояния.}

\glossitem{Механизм внимания (Attention mechanism)}{компонент нейронной сети, позволяющий модели фокусироваться на релевантных частях входных данных.}
